% =============================================================================
% CHAPTER 17: INTERVENTION FOUNDATIONS
% =============================================================================
% Type: C (Application Chapter)
% Version: 1.0 (January 2026)
% Status: Final
% Language: English
%
% Purpose: Establish the systematic framework for behavioral intervention design
% Primary Appendix: HHH (METHOD-TOOLKIT)
% Prerequisites: Chapter 13 (Journey), Chapter 14 (Segments), Chapter 15 (Willingness)
% Leads to: Chapter 18 (Journey-Integrated), Chapter 19 (Segment-Specific)
%
% =============================================================================

\chapter{Intervention Foundations}
\label{ch:intervention-foundations}

% -----------------------------------------------------------------------------
% METADATA BLOCK
% -----------------------------------------------------------------------------

\begin{tcolorbox}[colback=gray!5!white, colframe=gray!75!black, title=Chapter Metadata]
\small
\begin{tabular}{@{}ll@{}}
\textbf{Chapter Type:} & C (Application Chapter) \\
\textbf{Primary Appendix:} & HHH (METHOD-TOOLKIT: Intervention Design Toolkit) \\
\textbf{Prerequisites:} & Ch. 13 (Journey), Ch. 14 (Segments), Ch. 15 (Willingness) \\
\textbf{Leads to:} & Ch. 18 (Journey-Integrated), Ch. 19 (Segment-Specific) \\
\textbf{Language:} & English \\
\end{tabular}
\end{tcolorbox}

% -----------------------------------------------------------------------------
% APPENDIX REFERENCES BOX
% -----------------------------------------------------------------------------

\begin{tcolorbox}[colback=blue!5!white, colframe=blue!75!black,
    title=Appendix References]
\textbf{Primary:} Appendix TKT (METHOD-TOOLKIT) --- Complete 20-field intervention schema (Theorems HHH-T1 to HHH-T8)

\textbf{Supporting:}
\begin{itemize}[nosep]
\item Appendix CAL1 (METHOD-LLMMC-CAL): Calibration pipeline for R-Score decisions
\item Appendix UNMAPPED_STA (CORE-STAGE): Journey phases for F6 mapping
\item Appendix PAP1 (FORMAL-SEGMENT): Segment definitions for F9 targeting
\item Appendix UNMAPPED_WEN (CORE-WHEN): Context framework for F3 specification
\item Appendix UNMAPPED_PAP (FORMAL-EQUILIBRIA): Intervention equilibria theory
\end{itemize}

\textbf{Emergence Documentation:}
\begin{itemize}[nosep]
\item \texttt{docs/frameworks/core-framework-definition.yaml}: Single Source of Truth (SSOT) for 10C $\rightarrow$ HHH mapping
\item \texttt{data/gamma\_matrix\_complete.json}: $\gamma(\Psi \rightarrow \text{FEPSDE})$ feasibility matrix
\end{itemize}

\textbf{For Prioritization:} Appendix CAL1 provides the GO/PILOT/NO-GO decision framework
based on calibrated $\alpha$-parameters and R-Scores.
\end{tcolorbox}

% -----------------------------------------------------------------------------
% INTUITION BOX
% -----------------------------------------------------------------------------

\begin{tcolorbox}[colback=yellow!5!white, colframe=yellow!75!black,
    title=Intuition: Anna's Weight Loss Journey]

\textbf{Anna} wants to lose weight. She has tried many interventions:
\begin{itemize}[nosep]
\item Information campaigns (``Calories in, calories out'') --- didn't work
\item Gym membership (commitment device) --- worked for 2 months, then stopped
\item Fitbit feedback --- helpful but not sufficient
\item Social support group --- finally sustainable change
\end{itemize}

\textbf{Why did some interventions fail and others succeed?}

The answer lies in \textit{systematic intervention design}: matching the right intervention type to Anna's journey phase, segment characteristics, and context. A gym membership (commitment device) fails in the Awareness phase but succeeds in the Action phase. Social norms work for socially-oriented segments but not for autonomy-seekers.

This chapter provides the toolkit to design interventions systematically, not by trial and error.
\end{tcolorbox}

% -----------------------------------------------------------------------------
% CENTRAL QUESTION BOX
% -----------------------------------------------------------------------------

\begin{tcolorbox}[colback=red!5!white, colframe=red!75!black,
    title=Central Question]
\textbf{How do we systematically specify behavioral interventions to maximize effectiveness and enable organizational learning?}
\end{tcolorbox}

% -----------------------------------------------------------------------------
% METHODOLOGICAL NOTE: EXCLUSION PRINCIPLE
% -----------------------------------------------------------------------------
\begin{tcolorbox}[colback=red!5!white, colframe=red!75!black,
    title=\textbf{Methodological Note: The Exclusion Principle (Appendix UNMAPPED_FRM)}]
\textbf{Following Appendix UNMAPPED_FRM (LIT-FEHR-METHOD):}

Intervention design combines Journey phases (AW), Segments (BA), and Context ($\Psi$). Following the \textbf{Exclusion Principle}:
\begin{itemize}[nosep]
    \item \textbf{Start additive:} Test single-mechanism interventions first
    \item \textbf{Combination justified:} Only combine mechanisms if single mechanisms systematically fail
    \item \textbf{Scientific value:} Intervention combinations (``bundling'') gain value where they outperform the sum of parts
\end{itemize}

\textbf{Practical implication:} When designing multi-component interventions (e.g., information + incentive + commitment), first test whether simpler interventions work. Bundling must be empirically justified.
\end{tcolorbox}

% -----------------------------------------------------------------------------
% CHAPTER OVERVIEW
% -----------------------------------------------------------------------------

\begin{tcolorbox}[colback=green!5!white, colframe=green!75!black,
    title=Chapter Overview]

This chapter establishes the foundations for systematic intervention design:

\begin{enumerate}
\item \textbf{Section~\ref{sec:17-why}}: Why systematic intervention design matters
\item \textbf{Section~\ref{sec:17-types}}: The nine intervention types (from $W_{\text{eff}}$ decomposition)
\item \textbf{Section~\ref{sec:17-schema}}: The 20-field intervention schema
    \begin{itemize}[nosep]
    \item \textbf{Section~\ref{sec:17-emergence}}: Emergence from 10C Framework (14,400 $\rightarrow$ 2,000)
    \end{itemize}
\item \textbf{Section~\ref{sec:17-depths}}: Application depths (Light/Hybrid/Profound)
\item \textbf{Section~\ref{sec:17-axioms}}: The five intervention axioms (TKT-1 to TKT-5)
\item \textbf{Section~\ref{sec:17-transparency}}: Methodological transparency (Assumptions A1--A7)
\item \textbf{Section~\ref{sec:17-examples}}: Worked examples
\end{enumerate}
\end{tcolorbox}


%%%%%%%%%%%%%%%%%%%%%%%%%%%%%%%%%%%%%%%%%%%%%%%%%%%%%%%%%%%%%%%%%%%%%%%%%%%%%%%
\section{Why Systematic Intervention Design?}
\label{sec:17-why}
%%%%%%%%%%%%%%%%%%%%%%%%%%%%%%%%%%%%%%%%%%%%%%%%%%%%%%%%%%%%%%%%%%%%%%%%%%%%%%%

Behavioral interventions are everywhere. Governments use nudges to increase organ donation. Companies use defaults to boost retirement savings. Schools use feedback to improve student performance. Yet despite decades of research and billions invested, \textbf{most behavioral interventions underperform expectations}.

\subsection{The Problem with Ad-Hoc Design}

Consider the typical intervention design process:

\begin{enumerate}
\item Someone reads about a successful nudge (e.g., social norms for energy conservation)
\item They implement a similar intervention in their context
\item Results disappoint---maybe 5\% effect instead of the expected 20\%
\item The conclusion: ``Nudges don't work here''
\end{enumerate}

The real problem isn't that nudges don't work---it's that the intervention was designed \textit{ad hoc}, without systematic consideration of:

\begin{itemize}
\item \textbf{Context mismatch}: The original study was in California; your context is rural Switzerland
\item \textbf{Journey mismatch}: The intervention targeted Awareness, but your population is already in the Action phase
\item \textbf{Segment mismatch}: Social norms work for conformists, but your population includes many autonomy-seekers
\item \textbf{Combination conflict}: You added a financial incentive that crowded out the social motivation
\end{itemize}

\subsection{The Organizational Learning Problem}

Beyond individual intervention failures, organizations face a systemic problem: \textbf{they cannot learn from past interventions}.

\begin{itemize}
\item Intervention A was tried in 2019---but nobody documented why it was chosen or what happened
\item Intervention B worked in Department X---but Department Y doesn't know about it
\item A consultant recommended Intervention C---but left no documentation of the underlying logic
\end{itemize}

Without standardized documentation, every intervention design starts from zero. Organizational capability never accumulates.

\subsection{What Systematic Design Enables}

The EBF approach to intervention design solves these problems by providing:

\begin{enumerate}
\item \textbf{A common language}: The 20-field schema ensures everyone describes interventions the same way
\item \textbf{Explicit assumptions}: Fields F3 (Context), F6 (Journey), and F9 (Segment) force designers to state their assumptions
\item \textbf{Combinability analysis}: Field F10 (Complementarity) requires thinking about intervention interactions
\item \textbf{Organizational memory}: Standardized documentation enables learning across projects and time
\item \textbf{Predictive power}: With accumulated data, organizations can simulate intervention effects before deployment
\item \textbf{Prioritization decisions}: The R-Score framework (Appendix CAL1) enables GO/PILOT/NO-GO decisions based on calibrated fit parameters
\end{enumerate}

\begin{tcolorbox}[colback=blue!5!white, colframe=blue!75!black,
    title=The Promise of Systematic Design]
Organizations that adopt systematic intervention design can expect:
\begin{itemize}[nosep]
\item 30--50\% improvement in intervention effectiveness through better matching
\item 50\% reduction in failed interventions through explicit assumption testing
\item Cumulative organizational capability through standardized documentation
\item Predictable outcomes through simulation before deployment
\end{itemize}
\end{tcolorbox}


%%%%%%%%%%%%%%%%%%%%%%%%%%%%%%%%%%%%%%%%%%%%%%%%%%%%%%%%%%%%%%%%%%%%%%%%%%%%%%%
\section{The Nine Intervention Types}
\label{sec:17-types}
%%%%%%%%%%%%%%%%%%%%%%%%%%%%%%%%%%%%%%%%%%%%%%%%%%%%%%%%%%%%%%%%%%%%%%%%%%%%%%%

Every behavioral intervention operates through one or more \textit{mechanisms}---the psychological or structural pathway through which it affects behavior. The EBF framework identifies nine fundamental intervention types, each with a distinct mechanism.

\subsection{Type 1: Informative Interventions}

\textbf{Mechanism:} Knowledge provision---filling information gaps that prevent informed decision-making.

\textbf{Examples:}
\begin{itemize}[nosep]
\item Nutrition labels on food packaging
\item Health warning labels on cigarettes
\item Energy efficiency ratings on appliances
\item Financial literacy programs
\end{itemize}

\textbf{When it works:} When the target population lacks relevant information and would change behavior if informed. Classic examples include calorie labeling leading to healthier food choices.

\textbf{When it fails:} When awareness isn't the problem. Smokers know smoking is harmful---more information doesn't help. If the population is already in the Action phase, informative interventions are redundant.

\textbf{EBF connection:} Informative interventions primarily target the Awareness function $A(\cdot)$ from Appendix UNMAPPED_AWA, increasing $A_{\text{info}}$.

\subsection{Type 2: Normative Interventions}

\textbf{Mechanism:} Social norm activation---leveraging what others do (descriptive norms) or what others approve of (injunctive norms).

\textbf{Examples:}
\begin{itemize}[nosep]
\item ``9 out of 10 neighbors recycle'' (descriptive norm)
\item Hotel towel reuse signs referencing other guests
\item Tax compliance messages (``Most citizens pay on time'')
\item Voting reminders referencing community participation
\end{itemize}

\textbf{When it works:} When the target population is socially-oriented and the norm supports the desired behavior. \citet{allcott2011} showed 2--4\% energy reduction from social norm feedback.

\textbf{When it fails:} When the population is autonomy-oriented (reactance to perceived pressure) or when the actual norm contradicts the message (if only 40\% recycle, saying ``most people recycle'' backfires).

\textbf{EBF connection:} Normative interventions work through IDN (identity utility) and social preferences, activating $U^{IDN}$ from the FEPSDE framework.

\subsection{Type 3: Structural Interventions}

\textbf{Mechanism:} Choice architecture---changing the environment to make desired behaviors easier and undesired behaviors harder.

\textbf{Examples:}
\begin{itemize}[nosep]
\item Default enrollment in retirement savings \citep{madrian2001}
\item Opt-out organ donation systems
\item Smaller plates in cafeterias
\item Automatic escalation of savings rates
\item One-click checkout for desired products
\end{itemize}

\textbf{When it works:} When friction is a major barrier and the target population is ``passive''---willing to accept defaults. \citet{chetty2014} found 85\% of savers are passive.

\textbf{When it fails:} When the population has strong preferences that override defaults, or when the structural change is too visible and triggers reactance.

\textbf{EBF connection:} Structural interventions modify context $\Psi$, specifically $\Psi_I$ (institutional defaults) and $\Psi_P$ (physical environment).

\subsection{Type 4: Incentive-Based Interventions}

\textbf{Mechanism:} Financial rewards or penalties that change the cost-benefit calculation.

\textbf{Examples:}
\begin{itemize}[nosep]
\item Carbon taxes
\item Savings matching programs
\item Pay-for-performance in healthcare
\item Subsidies for electric vehicles
\item Penalties for non-compliance
\end{itemize}

\textbf{When it works:} When financial considerations are central to the decision and the incentive is large enough to matter.

\textbf{When it fails:} When intrinsic motivation is strong---adding payment can ``crowd out'' internal motivation. Classic example: paying children to read can reduce their interest in reading.

\textbf{EBF connection:} Incentive interventions directly modify $U^F$ (financial utility) in the FEPSDE framework. The crowding-out effect reflects negative complementarity with normative interventions ($\gamma < 0$).

\begin{tcolorbox}[colback=yellow!5!white, colframe=yellow!75!black,
    title=Warning: Crowding Out]
\textbf{Be careful combining incentives with social/identity interventions.}

When people do something because it's ``the right thing'' (intrinsic motivation), adding payment can signal that it's actually a transaction, not a moral act. The intrinsic motivation disappears, and when the payment stops, so does the behavior.

\textbf{Rule of thumb:} Don't pay people for things they would do anyway for social or identity reasons.
\end{tcolorbox}

\subsection{Type 5: Commitment Interventions}

\textbf{Mechanism:} Pre-commitment devices that lock in future behavior, overcoming present bias.

\textbf{Examples:}
\begin{itemize}[nosep]
\item Gym membership contracts with cancellation penalties
\item Savings commitment accounts \citep{ashraf2006}
\item Public goal announcements
\item Implementation intentions (``I will exercise at 7am on Monday'')
\item Deposit contracts for weight loss
\end{itemize}

\textbf{When it works:} When present bias ($\beta < 1$) is the problem and the person is a ``sophisticate'' who recognizes their self-control issues.

\textbf{When it fails:} When the person is a ``naif'' who doesn't recognize their present bias (they won't seek commitment). Also fails when commitment is too rigid and life circumstances change.

\textbf{EBF connection:} Commitment interventions address the Willingness function $WAX$ from Appendix UNMAPPED_REA, specifically the gap between $WAX_{\text{stated}}$ and $WAX_{\text{actual}}$.

\subsection{Type 6: Feedback Interventions}

\textbf{Mechanism:} Real-time information about behavior and its consequences, enabling self-correction.

\textbf{Examples:}
\begin{itemize}[nosep]
\item Smart meter energy displays
\item Fitness trackers and step counters
\item Spending alerts from banks
\item Speed display signs on roads
\item Real-time carbon footprint calculators
\end{itemize}

\textbf{When it works:} When people don't have accurate intuitions about their behavior and would adjust if they knew. Works best when feedback is immediate, salient, and actionable.

\textbf{When it fails:} When people already know their behavior (redundant information) or when feedback is too delayed or complex to act on.

\textbf{EBF connection:} Feedback interventions improve the accuracy of the Awareness function $A(\cdot)$ and support the Action phase of the behavioral change journey.

\subsection{Type 7: Identity-Based Interventions}

\textbf{Mechanism:} Activating or reshaping self-concept to align behavior with identity.

\textbf{Examples:}
\begin{itemize}[nosep]
\item ``Be a voter'' vs ``Vote'' framing \citep{bryan2011}
\item ``You are the kind of person who...'' messaging
\item Role model programs
\item Identity labels (``You are an environmentalist'')
\item Aspirational self framing
\end{itemize}

\textbf{When it works:} When identity is malleable and the target behavior can be framed as consistent with a desirable identity. Most effective in the Stabilization phase to lock in change.

\textbf{When it fails:} When identity is fixed or the proposed identity conflicts with existing self-concept.

\textbf{EBF connection:} Identity interventions work through $IDN^L$ (identity utility at level $L$), the most persistent form of behavioral change.

\subsection{Type 8: Temporal Interventions}

\textbf{Mechanism:} Using time strategically---deadlines, cooling-off periods, timing of prompts.

\textbf{Examples:}
\begin{itemize}[nosep]
\item Cooling-off periods for major purchases
\item Tax filing deadlines
\item ``Last chance'' urgency messaging
\item Optimal timing of reminders (payday for savings prompts)
\item Waiting periods for consequential decisions
\end{itemize}

\textbf{When it works:} When timing affects decision quality or when urgency/patience can be strategically induced.

\textbf{When it fails:} When artificial urgency is perceived as manipulation, or when cooling-off periods cause people to abandon beneficial decisions.

\textbf{EBF connection:} Temporal interventions modify $\Psi_T$ (temporal context) and interact with present bias ($\beta$) from the behavioral parameters.

\subsection{Type 9: Hybrid Interventions}

\textbf{Mechanism:} Combinations of multiple mechanisms in a single intervention.

\textbf{Examples:}
\begin{itemize}[nosep]
\item Gamified savings apps (feedback + incentives + social)
\item Weight Watchers (information + social norms + commitment)
\item Green energy programs (defaults + social norms + feedback)
\item Financial wellness programs (information + commitment + feedback)
\end{itemize}

\textbf{When it works:} When mechanisms are complementary ($\gamma > 0$) and target different barriers simultaneously.

\textbf{When it fails:} When mechanisms conflict ($\gamma < 0$)---e.g., social norms undermined by financial incentives.

\textbf{EBF connection:} Hybrid interventions require explicit complementarity analysis using Field F10 and the $\gamma_{ij}$ matrix.

\subsection{Summary: The Nine Types at a Glance}

\begin{table}[htbp]
\centering
\caption{The Nine Intervention Types}
\label{tab:17-types-summary}
\renewcommand{\arraystretch}{1.3}
\small
\begin{tabular}{@{}clll@{}}
\toprule
\textbf{\#} & \textbf{Type} & \textbf{Mechanism} & \textbf{Best For} \\
\midrule
1 & Informative & Knowledge provision & Awareness phase, information gaps \\
2 & Normative & Social norm activation & Social segments, conformity \\
3 & Structural & Choice architecture & Passive populations, friction \\
4 & Incentive-based & Financial rewards/penalties & Financial decisions \\
5 & Commitment & Pre-commitment devices & Present-biased sophisticates \\
6 & Feedback & Real-time behavior info & Action phase, self-monitoring \\
7 & Identity-based & Self-concept activation & Stabilization phase \\
8 & Temporal & Strategic use of time & Time-sensitive decisions \\
9 & Hybrid & Multi-mechanism & Complex behaviors \\
\bottomrule
\end{tabular}
\end{table}


%%%%%%%%%%%%%%%%%%%%%%%%%%%%%%%%%%%%%%%%%%%%%%%%%%%%%%%%%%%%%%%%%%%%%%%%%%%%%%%
\section{The 20-Field Intervention Schema}
\label{sec:17-schema}
%%%%%%%%%%%%%%%%%%%%%%%%%%%%%%%%%%%%%%%%%%%%%%%%%%%%%%%%%%%%%%%%%%%%%%%%%%%%%%%

The heart of systematic intervention design is the \textbf{20-field schema}---a standardized format for specifying any behavioral intervention. Just as accountants use standardized financial statements and engineers use standardized blueprints, behavioral designers should use standardized intervention specifications.

\subsection{Theoretical Foundation: Emergence from the 10C Framework}
\label{sec:17-emergence}

The 20-field schema is not arbitrary---it \textbf{emerges systematically} from the 10C CORE framework. This theoretical grounding ensures completeness and enables computational prediction.

\begin{tcolorbox}[colback=blue!5!white, colframe=blue!75!black,
    title=\textbf{The Emergence Architecture}]

\textbf{Step 1: 10C COREs define the classification dimensions}

Each CORE question generates a field dimension (see Appendix TKT, Theorems HHH-T1 to HHH-T7):

\begin{center}
\renewcommand{\arraystretch}{1.2}
\begin{tabular}{@{}llll@{}}
\toprule
\textbf{CORE} & \textbf{Question} & \textbf{Field} & \textbf{Categories} \\
\midrule
WHAT (C) & What is utility? & F5 (FEPSDE) & 6 dimensions \\
WHO (AAA) & Who has utility? & F4 (Target) & 4 levels \\
HOW (B) & How interact? & F10 (Complementarity) & $\gamma$ values \\
WHEN (V) & When context? & F3 (Context) & 8 $\Psi$-dimensions \\
STAGE (AW) & Where in journey? & F6 (Phase) & 5 phases \\
AWARE (AU) & How aware? & (F21 proposed) & 3 regimes \\
HIERARCHY (HI) & Which decision level? & F12 (Scope) & 5 scopes \\
\bottomrule
\end{tabular}
\end{center}

\textbf{Step 2: W\textsubscript{eff} decomposition defines intervention types}

From the Effective Willingness function (Theorem HHH-T1):
\begin{equation}
W_{\text{eff}} = W_{\text{base}} \cdot A(\cdot) \cdot (1 + \kappa_{AWX}) \cdot (1 + \kappa_{KON}) \cdot (1 + \kappa_{JNY})
\end{equation}

Each component that can be targeted yields an intervention type $\Rightarrow$ \textbf{8 primitive types} (Field F2).

\end{tcolorbox}

\begin{tcolorbox}[colback=orange!5!white, colframe=orange!75!black,
    title=\textbf{The 14,400 Archetype Space}]

\textbf{Theorem HHH-T7 (Complete Classification):}

Combining the emergent dimensions yields a 6-dimensional classification space:
\begin{equation}
\mathcal{I} \mapsto (\underbrace{\text{Type}}_{8}, \underbrace{\text{Scope}}_{5}, \underbrace{\text{Level}}_{4}, \underbrace{\text{FEPSDE}}_{6}, \underbrace{\text{Phase}}_{5}, \underbrace{\text{Awareness}}_{3})
\end{equation}

\textbf{Theoretical space:} $8 \times 5 \times 4 \times 6 \times 5 \times 3 = \mathbf{14{,}400}$ intervention archetypes.

\textbf{The challenge:} Most of these combinations are infeasible or redundant. We need a principled reduction method.

\end{tcolorbox}

\begin{tcolorbox}[colback=green!5!white, colframe=green!75!black,
    title=\textbf{The Elegant Reduction: $\gamma(\Psi)$ Feasibility Filter}]

\textbf{Theorem HHH-T8 (Feasibility Constraints):}

An archetype $\vec{a}$ is feasible if and only if it passes four constraint classes:
\begin{equation}
\text{Feasible}(\vec{a}) \iff C_{\text{affinity}} \wedge C_{\text{logical}} \wedge C_{\text{scale}} \wedge C_{\text{empirical}}
\end{equation}

\textbf{The key insight:} Context-dependent complementarity $\gamma(\Psi)$ provides a universal filter:
\begin{equation}
\boxed{\text{Feasible}(\mathcal{I} | \Psi) \iff \gamma(\Psi \rightarrow \text{target\_FEPSDE}(\mathcal{I})) \geq \theta}
\end{equation}

\textbf{Result:} $14{,}400 \rightarrow \approx 2{,}000$ feasible archetypes (context-dependent).

\textbf{Data source:} The $\gamma(\Psi \rightarrow \text{FEPSDE})$ matrix is estimated via LLM Monte Carlo (Appendix CAL1) and stored in \texttt{data/gamma\_matrix\_complete.json}.

\end{tcolorbox}

\begin{tcolorbox}[colback=gray!5!white, colframe=gray!75!black,
    title=Implication for Practice]
The 20-field schema captures:
\begin{itemize}[nosep]
    \item \textbf{6 emergent fields} (F2, F4, F5, F6, F12, +F21) --- derived from 10C via theorems
    \item \textbf{1 context field} (F3) --- maps directly to $\Psi$
    \item \textbf{2 interface fields} (F9, F10) --- links to Segments (BA) and Complementarity (B)
    \item \textbf{11 operational fields} (F1, F7, F8, F11, F13--F20) --- implementation details
\end{itemize}

\textbf{Why this matters:} Because the core fields emerge from theory, we can \textit{predict} intervention effectiveness before deployment. The operational fields ensure practical implementability.

\textbf{Single Source of Truth:} The complete mapping is documented in \texttt{docs/frameworks/core-framework-definition.yaml} (Section: derived\_structures).
\end{tcolorbox}

\subsection{Why 20 Fields?}

The schema captures everything needed to:
\begin{itemize}
\item \textbf{Describe} the intervention (F1, F2, F19)
\item \textbf{Position} it in context (F3), journey (F6), and segment (F9)
\item \textbf{Target} specific utility dimensions (F4, F5)
\item \textbf{Design} the delivery (F7, F8, F12, F14, F15)
\item \textbf{Anticipate} interactions (F10, F11)
\item \textbf{Evaluate} effectiveness (F13, F17, F20)
\item \textbf{Assign} responsibility (F16, F18)
\end{itemize}

\subsection{The Core Fields: Light Mode (F1--F6)}

Every intervention must specify at least these six fields---the \textit{minimum viable specification}:

\begin{tcolorbox}[colback=green!5!white, colframe=green!75!black,
    title=Light Mode: The Essential Six Fields]

\textbf{F1: Title/Label}\\
A clear, descriptive name. Not ``Intervention A'' but ``Auto-Enrollment Retirement Default''.

\textbf{F2: Intervention Type}\\
One of the nine types (Section~\ref{sec:17-types}). If hybrid, specify which types are combined.

\textbf{F3: Context Level}\\
Where is this implemented?
\begin{itemize}[nosep]
\item State (national/regional government)
\item Organization (company, NGO, institution)
\item Family/Household
\item Individual
\item Digital (platform, app)
\end{itemize}

\textbf{F4: Target Structure}\\
Which utility type does this target?
\begin{itemize}[nosep]
\item IND (Individual utility---personal benefit)
\item IDN (Identity utility---who I am)
\item Collective (Group/social benefit)
\end{itemize}

\textbf{F5: FEPSDE Dimension}\\
Which welfare dimension(s)?
\begin{itemize}[nosep]
\item Financial, Emotional, Physical, Social, Digital, Ecological
\end{itemize}

\textbf{F6: Journey Phase}\\
What phase is the target population in?
\begin{itemize}[nosep]
\item Awareness, Trigger, Action, Maintenance, Stabilization
\end{itemize}

\end{tcolorbox}

\textbf{Time required:} 10 minutes for a skilled practitioner. This is the ``napkin test''---if you can't fill out F1--F6 quickly, you don't understand your intervention well enough.

\subsection{Extended Fields: Hybrid Mode (F7--F12)}

For serious intervention design, add these six fields:

\begin{table}[htbp]
\centering
\caption{Hybrid Mode Fields (F7--F12)}
\label{tab:17-hybrid-fields}
\renewcommand{\arraystretch}{1.3}
\small
\begin{tabular}{@{}clp{8cm}@{}}
\toprule
\textbf{Field} & \textbf{Name} & \textbf{Description} \\
\midrule
F7 & Framing Logic & How is the intervention framed? Positive (gain), Negative (loss), Social (others), Neutral \\
F8 & Autonomy Level & Voluntary choice, Default (opt-out), Mandate, Prohibition \\
F9 & Target Segment & Which behavioral segment? (Reference to Chapter 14 segment IDs) \\
F10 & Complementarity & Links to other interventions with $\gamma$ estimates \\
F11 & Side Effect Risk & Low/Medium/High + explanation of potential negative effects \\
F12 & Time Horizon & Instant, Short-term (days), Medium (months), Long-term (years) \\
\bottomrule
\end{tabular}
\end{table}

\textbf{Time required:} 30 minutes. This is the ``project proposal'' level---enough detail for a steering committee decision.

\subsection{Full Specification: Profound Mode (F13--F20)}

For research-grade documentation or high-stakes interventions:

\begin{table}[htbp]
\centering
\caption{Profound Mode Fields (F13--F20)}
\label{tab:17-profound-fields}
\renewcommand{\arraystretch}{1.3}
\small
\begin{tabular}{@{}clp{8cm}@{}}
\toprule
\textbf{Field} & \textbf{Name} & \textbf{Description} \\
\midrule
F13 & Evaluation Plan & Yes/No + method (RCT, quasi-experiment, survey, KPIs) \\
F14 & Path Function & Role in journey: Door-opener, Amplifier, Escalator, Stabilizer \\
F15 & Repetition Required & One-time, Cyclic (weekly/monthly), Continuous \\
F16 & Responsible Role & Who implements? (HR, Policy team, External vendor, etc.) \\
F17 & Source/Evidence & Literature references, pilot results, expert judgment \\
F18 & System Link/ID & Internal tracking ID, API reference, database key \\
F19 & Description & Full narrative explanation of mechanism and rationale \\
F20 & Status & Planned, Piloted, Active, Scaled, Retired \\
\bottomrule
\end{tabular}
\end{table}

\textbf{Time required:} 60 minutes. This is the ``academic paper'' level---complete documentation for organizational learning and external communication.


%%%%%%%%%%%%%%%%%%%%%%%%%%%%%%%%%%%%%%%%%%%%%%%%%%%%%%%%%%%%%%%%%%%%%%%%%%%%%%%
\section{Application Depths}
\label{sec:17-depths}
%%%%%%%%%%%%%%%%%%%%%%%%%%%%%%%%%%%%%%%%%%%%%%%%%%%%%%%%%%%%%%%%%%%%%%%%%%%%%%%

Not every intervention needs full specification. The schema supports three \textit{application depths}, allowing proportionate effort:

\subsection{Light Mode: Rapid Screening}

\textbf{Fields:} F1--F6\\
\textbf{Time:} 10 minutes\\
\textbf{Use when:}
\begin{itemize}[nosep]
\item Brainstorming intervention options
\item Initial screening of ideas
\item Quick documentation of existing interventions
\item Communication with non-specialists
\end{itemize}

\textbf{Output:} A one-page intervention card suitable for portfolio overview.

\subsection{Hybrid Mode: Standard Design}

\textbf{Fields:} F1--F12\\
\textbf{Time:} 30 minutes\\
\textbf{Use when:}
\begin{itemize}[nosep]
\item Designing new interventions for implementation
\item Preparing project proposals
\item Comparing intervention alternatives
\item Building intervention portfolios
\end{itemize}

\textbf{Output:} A complete intervention specification suitable for implementation decision.

\subsection{Profound Mode: Research Grade}

\textbf{Fields:} F1--F18 (F19--F20 optional)\\
\textbf{Time:} 60 minutes\\
\textbf{Use when:}
\begin{itemize}[nosep]
\item High-stakes or high-cost interventions
\item Research publications
\item Regulatory submissions
\item Building organizational knowledge base
\end{itemize}

\textbf{Output:} Complete documentation suitable for academic publication or regulatory review.

\begin{tcolorbox}[colback=blue!5!white, colframe=blue!75!black,
    title=Progression Principle]
\textbf{Start Light, deepen as needed.}

Most intervention ideas should be screened at Light Mode first. Only promising candidates advance to Hybrid Mode for serious design work. Profound Mode is reserved for interventions that will actually be implemented at scale or published.

This prevents over-investment in early-stage ideas while ensuring thorough documentation for important interventions.
\end{tcolorbox}


%%%%%%%%%%%%%%%%%%%%%%%%%%%%%%%%%%%%%%%%%%%%%%%%%%%%%%%%%%%%%%%%%%%%%%%%%%%%%%%
\section{The Five Intervention Axioms}
\label{sec:17-axioms}
%%%%%%%%%%%%%%%%%%%%%%%%%%%%%%%%%%%%%%%%%%%%%%%%%%%%%%%%%%%%%%%%%%%%%%%%%%%%%%%

Five axioms govern effective intervention design. Violating any of them predicts intervention failure.

\subsection{Axiom TKT-1: Completeness}

\begin{tcolorbox}[colback=gray!5!white, colframe=gray!75!black,
    title=Axiom TKT-1: Intervention Completeness]
\textbf{An intervention can only be evaluated or deployed if all required fields for its application depth are specified.}

\textbf{Implication:} Don't implement interventions with missing specifications. If you can't fill out the Light Mode fields, you don't understand your intervention well enough to deploy it.
\end{tcolorbox}

\subsection{Axiom TKT-2: Context Dependency}

\begin{tcolorbox}[colback=gray!5!white, colframe=gray!75!black,
    title=Axiom TKT-2: Context Dependency]
\textbf{No intervention has context-independent effectiveness. The same intervention works differently in different contexts.}

\textbf{Implication:} Always specify F3 (Context Level). Never assume that because an intervention worked in California, it will work in Switzerland. Context includes culture, institutions, physical environment, and temporal factors.
\end{tcolorbox}

\subsection{Axiom TKT-3: Journey Phase Alignment}

\begin{tcolorbox}[colback=gray!5!white, colframe=gray!75!black,
    title=Axiom TKT-3: Journey Phase Alignment]
\textbf{Intervention effectiveness is maximized when the intervention type matches the target population's actual journey phase.}

Misaligned interventions achieve less than 30\% of their potential effectiveness.

\textbf{Implication:} Diagnose the journey phase \textit{before} selecting intervention type. An informative intervention fails if the population is already aware. A commitment device fails if the population hasn't formed intention yet.
\end{tcolorbox}

\begin{tcolorbox}[colback=yellow!5!white, colframe=yellow!75!black,
    title=Example: Anna's Misaligned Intervention]
Remember Anna from the chapter opening? Her gym membership (commitment device) failed because she was still in the \textbf{Awareness phase}---she knew she should exercise but hadn't formed a concrete intention. A commitment device is for the \textbf{Trigger phase}, when intention exists but action is delayed.

The social support group worked because it addressed the right phase (Maintenance/Stabilization) and matched her segment (socially-oriented).
\end{tcolorbox}

\subsection{Axiom TKT-4: Complementarity Structure}

\begin{tcolorbox}[colback=gray!5!white, colframe=gray!75!black,
    title=Axiom TKT-4: Complementarity Structure]
\textbf{Intervention pairs can reinforce each other ($\gamma > 0$), have independent effects ($\gamma = 0$), or undermine each other ($\gamma < 0$).}

The combined effect is:
\begin{equation}
E(\mathcal{I}_i + \mathcal{I}_j) = E(\mathcal{I}_i) + E(\mathcal{I}_j) + \gamma_{ij} \cdot \sqrt{E(\mathcal{I}_i) \cdot E(\mathcal{I}_j)}
\end{equation}

\textbf{Implication:} Before combining interventions, check complementarity. Avoid combinations with $\gamma < 0$ (e.g., financial incentives undermining social motivation).
\end{tcolorbox}

\subsection{Axiom TKT-5: Segment Specificity}

\begin{tcolorbox}[colback=gray!5!white, colframe=gray!75!black,
    title=Axiom TKT-5: Segment Specificity]
\textbf{The same intervention can have dramatically different effects on different behavioral segments---including opposite effects.}

Effectiveness varies by segment multiplier $\sigma_s \in [0, 2]$ or even negative (reactance).

\textbf{Implication:} Always consider who you're targeting. A normative intervention might boost behavior in socially-oriented segments ($\sigma = 1.6$) while triggering reactance in autonomy-oriented segments ($\sigma = -0.3$).
\end{tcolorbox}


%%%%%%%%%%%%%%%%%%%%%%%%%%%%%%%%%%%%%%%%%%%%%%%%%%%%%%%%%%%%%%%%%%%%%%%%%%%%%%%
\section{Methodological Transparency: Emergence vs. Assumptions}
\label{sec:17-transparency}
%%%%%%%%%%%%%%%%%%%%%%%%%%%%%%%%%%%%%%%%%%%%%%%%%%%%%%%%%%%%%%%%%%%%%%%%%%%%%%%

Scientific integrity requires explicit documentation of what is \textit{derived from theory} versus what is \textit{assumed}. This section provides full transparency about the reduction pipeline from 14,400 theoretical archetypes to $\approx$2,000 feasible interventions.

\subsection{The Reduction Pipeline}

\begin{table}[htbp]
\centering
\caption{Reduction Pipeline: Emergence vs. Assumptions}
\label{tab:17-reduction-pipeline}
\renewcommand{\arraystretch}{1.3}
\small
\begin{tabular}{@{}clccp{4.5cm}@{}}
\toprule
\textbf{Stage} & \textbf{Description} & \textbf{Output} & \textbf{Status} & \textbf{Assumptions} \\
\midrule
0 & 10C CORE Framework & 9 dimensions & \textcolor{green!60!black}{\textbf{Emergent}} & None (50+ years research) \\
1 & 6D Classification & 14,400 & Mixed & A1: Cardinalities \\
2 & Affinity Constraints & $\approx$8,000 & \textcolor{orange}{\textbf{Assumed}} & A2: Matrix, A3: $\theta_\alpha$ \\
3 & Logical Constraints & $\approx$5,000 & \textcolor{green!60!black}{\textbf{Emergent}} & None (logical necessity) \\
4 & Scale Constraints & $\approx$3,500 & \textcolor{orange}{\textbf{Assumed}} & A4: Coherence $\leq 2$ \\
5 & Empirical Constraints & $\approx$2,500 & \textcolor{red!60!black}{\textbf{Assumed}} & A5: Evidence required \\
6 & $\gamma(\Psi)$ Filter & $\approx$2,000 & \textcolor{orange}{\textbf{Assumed}} & A6: Matrix, A7: $\theta$ \\
\bottomrule
\end{tabular}
\end{table}

\subsection{Explicit Assumption Register}

\begin{tcolorbox}[colback=yellow!5!white, colframe=yellow!75!black,
    title=\textbf{Assumption A1: Dimension Cardinalities}]
\textbf{Assumption:} The categorical discretization of continuous dimensions.

\begin{itemize}[nosep]
\item Journey Phases: 5 (could be 3 or 7)
\item Awareness Regimes: 3 (could be continuous $A \in [0,1]$)
\item Temporal Scopes: 5 (could be 4 or 6)
\end{itemize}

\textbf{Justification:} Based on empirical clustering in behavioral research (e.g., Transtheoretical Model stages).

\textbf{Risk:} Granularity loss---nuanced differences within categories are obscured.

\textbf{Mitigation:} Cardinalities can be increased if needed; continuous versions available in Appendix TKT.
\end{tcolorbox}

\begin{tcolorbox}[colback=yellow!5!white, colframe=yellow!75!black,
    title=\textbf{Assumptions A2--A3: Affinity Constraints}]
\textbf{Assumption A2:} Type-Phase and Type-Awareness affinity matrices exist with values $\alpha_{t,\varphi} \in [0,1]$.

\textbf{Assumption A3:} Threshold $\theta_\alpha = 0.4$ for feasibility.

\textbf{Justification:} Matrix values derived from meta-analyses of intervention effectiveness by phase. Threshold chosen to exclude weak but not marginal combinations.

\textbf{Risk:} \textbf{False negatives}---valid interventions may be excluded if affinity is underestimated.

\textbf{Mitigation:}
\begin{itemize}[nosep]
\item Lower $\theta_\alpha$ for exploratory design ($\theta_\alpha = 0.2$)
\item Empirically recalibrate affinity matrices with new evidence
\item Flag borderline cases ($0.3 \leq \alpha < 0.5$) for pilot testing
\end{itemize}
\end{tcolorbox}

\begin{tcolorbox}[colback=yellow!5!white, colframe=yellow!75!black,
    title=\textbf{Assumption A4: Scale Coherence}]
\textbf{Assumption:} $|\text{scope\_rank} - \text{level\_rank}| \leq 2$

\textbf{Justification:} Individual-level interventions cannot plausibly have system-wide temporal scope; societal interventions cannot be instant.

\textbf{Risk:} Some innovative cross-scale interventions may be excluded.

\textbf{Mitigation:} Relax constraint to $\leq 3$ for innovative search mode.
\end{tcolorbox}

\begin{tcolorbox}[colback=red!5!white, colframe=red!75!black,
    title=\textbf{Assumption A5: Empirical Evidence Requirement} $\triangleright$ \textit{Highest Risk}]
\textbf{Assumption:} A (Type, FEPSDE) pairing requires literature support.

\textbf{Critical Problem:}
\begin{equation}
\text{Absence of Evidence} \neq \text{Evidence of Absence}
\end{equation}

Novel combinations that have never been tested are excluded---but they might work.

\textbf{Example:} ``Identity-based Ecological Interventions''---limited literature, but potentially effective for climate behavior.

\textbf{Risk:} \textbf{Highest false negative risk}---innovation is systematically suppressed.

\textbf{Mitigation: Three-Tier Archetype Classification}

Aligned with the Epistemic Tag System (see \texttt{quality/instruments/epistemic\_tags.md}):

\begin{itemize}[nosep]
\item \textcolor{green!60!black}{\textbf{EMP-Archetype}}: Empirically tested combination ($\checkmark$ scale)
\item \textcolor{blue!60!black}{\textbf{THR-Archetype}}: Theoretically plausible, not yet tested ($\checkmark$ pilot)
\item \textcolor{orange}{\textbf{EXP-Archetype}}: Novel combination, no evidence ($\checkmark$ experiment)
\end{itemize}

\medskip
\textbf{Relationship to Parameter-Level Epistemic Tags:}

\begin{center}
\small
\renewcommand{\arraystretch}{1.2}
\begin{tabular}{@{}lll@{}}
\toprule
\textbf{Level} & \textbf{Parameter Tag} & \textbf{Archetype Status} \\
\midrule
Validated & EMP (empirically validated) & EMP-Archetype \\
Derived & THR (theoretically derived) & THR-Archetype \\
Estimated & LLM (LLMMC-estimated) & --- \\
Exploratory & ILL/HYP (illustrative) & EXP-Archetype \\
\bottomrule
\end{tabular}
\end{center}

\textbf{Key distinction:} An archetype can use EMP-parameters but still be THR-Archetype if the \textit{combination} was never tested.

\textbf{Recommendation:} Do NOT use A5 as hard filter. Use archetype status for prioritization:
\begin{itemize}[nosep]
\item EMP-Archetype $\rightarrow$ Scale confidently
\item THR-Archetype $\rightarrow$ Pilot with measurement
\item EXP-Archetype $\rightarrow$ Small-scale experiment first
\end{itemize}
\end{tcolorbox}

\begin{tcolorbox}[colback=yellow!5!white, colframe=yellow!75!black,
    title=\textbf{Assumptions A6--A7: $\gamma(\Psi)$ Context Filter}]
\textbf{Assumption A6:} The $\gamma(\Psi \rightarrow \text{FEPSDE})$ matrix (48 coefficients) accurately captures context-dependent complementarity.

\textbf{Assumption A7:} Feasibility threshold $\theta = 0.3$ is appropriate.

\textbf{Justification:} Matrix estimated via LLM Monte Carlo (Appendix CAL1) with uncertainty quantification.

\textbf{Risk:} LLMMC estimates may be biased or have high variance.

\textbf{Mitigation:}
\begin{itemize}[nosep]
\item Carry forward confidence intervals from LLMMC
\item Empirically recalibrate with field data
\item Use conservative threshold ($\theta = 0.4$) when uncertainty is high
\end{itemize}
\end{tcolorbox}

\subsection{Risk Assessment: What Solutions Might We Lose?}

\begin{table}[htbp]
\centering
\caption{Risk Assessment for Each Reduction Stage}
\label{tab:17-risk-assessment}
\renewcommand{\arraystretch}{1.3}
\small
\begin{tabular}{@{}clccp{4cm}@{}}
\toprule
\textbf{Stage} & \textbf{Constraint} & \textbf{Risk Level} & \textbf{Type} & \textbf{Lost Solutions} \\
\midrule
1 & Cardinalities & Low & Granularity & Nuanced sub-phases \\
2 & Affinity & Medium & False Negative & Low-affinity but valid \\
3 & Logical & \textcolor{green!60!black}{None} & --- & Impossible combinations \\
4 & Scale & Low--Medium & False Negative & Cross-scale innovations \\
5 & Empirical & \textcolor{red!60!black}{\textbf{High}} & False Negative & Novel untested combinations \\
6 & $\gamma(\Psi)$ & Medium & Estimation Error & Context misclassification \\
\bottomrule
\end{tabular}
\end{table}

\begin{tcolorbox}[colback=blue!5!white, colframe=blue!75!black,
    title=Key Insight: The Exploration-Exploitation Tradeoff]

The reduction pipeline trades \textbf{exploration} (finding novel solutions) for \textbf{exploitation} (focusing on validated solutions).

\textbf{Default mode:} Exploitation-focused (all constraints active) $\rightarrow$ $\approx$2,000 archetypes

\textbf{Exploration mode:} Relaxed constraints $\rightarrow$ $\approx$5,000+ archetypes
\begin{itemize}[nosep]
\item Lower $\theta_\alpha$ (A3)
\item Relax scale coherence (A4)
\item Disable empirical filter (A5) $\rightarrow$ use EMP/THR/EXP-Archetype classification instead
\item Widen $\gamma$ threshold (A7)
\end{itemize}

\textbf{Recommendation:} Use exploration mode for innovation projects; exploitation mode for scaling proven interventions.
\end{tcolorbox}

\subsection{Validation Pathway}

For any intervention that emerges from the pipeline:

\begin{enumerate}
\item \textbf{Stage 1 (Design):} Use 20-field schema with appropriate mode (Light/Hybrid/Profound)
\item \textbf{Stage 2 (Simulate):} Estimate R-Score using Appendix CAL1 calibration pipeline
\item \textbf{Stage 3 (Pilot):} Test with small sample before scaling
\item \textbf{Stage 4 (Validate):} Compare predicted vs. actual effectiveness
\item \textbf{Stage 5 (Recalibrate):} Update $\gamma$ matrix and affinity values with empirical data
\end{enumerate}

This creates a \textbf{learning loop}: each implemented intervention improves future predictions.


%%%%%%%%%%%%%%%%%%%%%%%%%%%%%%%%%%%%%%%%%%%%%%%%%%%%%%%%%%%%%%%%%%%%%%%%%%%%%%%
\section{Worked Examples}
\label{sec:17-examples}
%%%%%%%%%%%%%%%%%%%%%%%%%%%%%%%%%%%%%%%%%%%%%%%%%%%%%%%%%%%%%%%%%%%%%%%%%%%%%%%

\subsection{Example 1: Retirement Savings Auto-Enrollment}

\textbf{Background:} A pension fund wants to increase retirement savings participation among employees aged 25--35. Current voluntary participation rate is 45\%.

\textbf{Light Mode Specification:}

\begin{table}[htbp]
\centering
\small
\renewcommand{\arraystretch}{1.3}
\begin{tabular}{@{}llp{6cm}@{}}
\toprule
\textbf{Field} & \textbf{Value} & \textbf{Rationale} \\
\midrule
F1: Title & Auto-Enrollment Default & Clear, descriptive \\
F2: Type & Structural (Type 3) & Choice architecture \\
F3: Context & Organizational & Employer implementation \\
F4: Target & IND (Future Self) & Long-term individual benefit \\
F5: FEPSDE & Financial & Savings behavior \\
F6: Journey & Trigger $\rightarrow$ Action & Converting intention to action \\
\bottomrule
\end{tabular}
\end{table}

\textbf{Result:} Participation increases from 45\% to 86\% \citep{madrian2001}---nearly doubling with minimal cost.

\textbf{Hybrid Extension:}
\begin{itemize}
\item F7 (Framing): Neutral---default presented as standard
\item F8 (Autonomy): Default with easy opt-out
\item F10 (Complementarity): High synergy with auto-escalation ($\gamma = 0.6$)
\item F11 (Side Effects): Low---preserves choice
\item F12 (Horizon): Long-term (career-long effects)
\end{itemize}

\textbf{Key insight:} This intervention works because 85\% of savers are ``passive''---they accept defaults. For the 15\% who are active, opt-out preserves choice.

\subsection{Example 2: Energy Conservation Portfolio}

\textbf{Background:} A municipality wants to reduce household energy consumption by 15\% over 3 years.

\textbf{Portfolio Design:} Rather than a single intervention, the municipality designs a \textit{coherent portfolio} covering the full journey:

\begin{table}[htbp]
\centering
\small
\renewcommand{\arraystretch}{1.3}
\begin{tabular}{@{}lllccl@{}}
\toprule
\textbf{Intervention} & \textbf{Type} & \textbf{Journey} & \textbf{F14} & \textbf{$\gamma$} \\
\midrule
Social norm mailers & Normative & Awareness & Door-opener & -- \\
Smart meter feedback & Feedback & Action & Amplifier & 0.5 \\
Commitment cards & Commitment & Maintenance & Stabilizer & 0.4 \\
\bottomrule
\end{tabular}
\end{table}

\textbf{Complementarity check:}
\begin{itemize}
\item Social norms + Feedback: $\gamma = 0.5$ (synergy---norms motivate, feedback enables)
\item Feedback + Commitment: $\gamma = 0.4$ (synergy---commitment locks in feedback gains)
\item Social norms + Commitment: $\gamma = 0.3$ (moderate synergy)
\end{itemize}

All $\gamma > 0$---the portfolio is coherent.

\textbf{Result:} Combined effect of 12--18\% reduction \citep{allcott2011}, exceeding what any single intervention would achieve.

\textbf{Key insight:} The portfolio works because each intervention addresses a different journey phase, and all three reinforce each other. A single intervention would have achieved perhaps 5--8\%.

\subsection{Example 3: A Failed Intervention (What Not to Do)}

\textbf{Background:} A company tries to increase employee volunteering by combining a social norm message (``70\% of your colleagues volunteer'') with a financial bonus (``\$50 for every volunteer day'').

\textbf{Light Mode Analysis:}
\begin{itemize}
\item F2: Hybrid (Normative + Incentive)
\item F10: Complementarity = ?
\end{itemize}

\textbf{The problem:} Normative and incentive-based interventions often have \textit{negative} complementarity ($\gamma < 0$). The financial payment signals that volunteering is a transaction, not a moral act. This ``crowds out'' the social motivation.

\textbf{Result:} Volunteering actually \textit{decreases} among socially-motivated employees. The intervention backfires.

\textbf{Key insight:} Always check complementarity before combining interventions. Some combinations that seem intuitively good are actually harmful.


%%%%%%%%%%%%%%%%%%%%%%%%%%%%%%%%%%%%%%%%%%%%%%%%%%%%%%%%%%%%%%%%%%%%%%%%%%%%%%%
\section{Summary}
\label{sec:17-summary}
%%%%%%%%%%%%%%%%%%%%%%%%%%%%%%%%%%%%%%%%%%%%%%%%%%%%%%%%%%%%%%%%%%%%%%%%%%%%%%%

\begin{tcolorbox}[colback=green!10!white, colframe=green!75!black,
    title=Chapter 17 Summary]

\textbf{Core Message:} Behavioral interventions require systematic design, not trial and error.

\textbf{The Nine Intervention Types:}
\begin{enumerate}[nosep]
\item Informative (knowledge)
\item Normative (social norms)
\item Structural (choice architecture)
\item Incentive-based (financial)
\item Commitment (pre-commitment)
\item Feedback (real-time info)
\item Identity-based (self-concept)
\item Temporal (timing)
\item Hybrid (combinations)
\end{enumerate}

\textbf{The 20-Field Schema:}
\begin{itemize}[nosep]
\item Light Mode (F1--F6): 10 minutes, minimum viable specification
\item Hybrid Mode (F1--F12): 30 minutes, standard design
\item Profound Mode (F1--F18): 60 minutes, research-grade
\end{itemize}

\textbf{The Five Axioms:}
\begin{itemize}[nosep]
\item TKT-1: Completeness---specify all required fields
\item TKT-2: Context Dependency---no universal interventions
\item TKT-3: Journey Alignment---match type to phase
\item TKT-4: Complementarity---check $\gamma$ before combining
\item TKT-5: Segment Specificity---same intervention, different effects
\end{itemize}

\textbf{Key Takeaway:} Use the schema. Check the axioms. Design portfolios, not isolated interventions.

\end{tcolorbox}


% -----------------------------------------------------------------------------
% READING PATH BOX
% -----------------------------------------------------------------------------

\begin{tcolorbox}[colback=purple!5!white, colframe=purple!75!black,
    title=Reading Path]

\textbf{You have completed:} Chapter 17 (Intervention Foundations)

\textbf{You now understand:}
\begin{itemize}[nosep]
\item The nine intervention types and their mechanisms
\item The 20-field schema for intervention specification
\item The three application depths (Light/Hybrid/Profound)
\item The five axioms governing effective design
\end{itemize}

\textbf{Next steps:}
\begin{itemize}[nosep]
\item \textbf{Chapter 18}: Journey-Integrated Interventions --- Which intervention for which journey phase?
\item \textbf{Chapter 19}: Segment-Specific Interventions --- How to target by behavioral type
\item \textbf{Chapter 20}: Intervention Portfolios --- How to combine interventions optimally
\end{itemize}

\textbf{For formal details:}
\begin{itemize}[nosep]
\item Appendix TKT (METHOD-TOOLKIT): Complete 20-field intervention schema
\item Appendix CAL1 (METHOD-LLMMC-CAL): R-Score calibration and GO/PILOT/NO-GO decisions
\end{itemize}
\end{tcolorbox}


% =============================================================================
% END OF CHAPTER 17
% =============================================================================
